%!TEX root = TheoretischePhysikC.tex
% Kern von content/c03-hilbertraum, c03.s13 bis c03.s15. Quellen: docs/plan/script-v0.3/sync-ch5.md

\section{Vorlesung: Spur, Diagonalisierung, Funktionalkalkül}

Mit dieser Vorlesung ist das mathematische Werkzeug der Quantenmechanik vollständig. Die Positions-Basis $\ket{x}$, mit der diese Vorlesung früher begann, kommt mit dem Teilchen auf der Geraden in Vorlesung~7.

\subsection{Spur}
Aus der Matrizenrechnung kennen Sie die Summe der Diagonaleinträge. Sie kommt in der Quantenmechanik ständig vor.
\begin{definition}
	Sei $\{\ket{x}\}_{x\in X}$ eine Orthonormalbasis von $\Hil$ mit endlicher Indexmenge; $x$ ist hier nur ein Index. Die \emph{Spur} eines Operators $\op{A}$ ist
	\begin{equation}
		\tr(\op{A}) = \sum_{x\in X} \melem{x}{\op{A}}{x}.
	\end{equation}
\end{definition}
Die Zahl hängt nicht von der Orthonormalbasis ab: Schiebt man die Vollständigkeitsrelation $\mathbb{I} = \sum_y \ketbra{y}{y}$ einer zweiten Basis ein, so wird $\sum_x \melem{x}{\op{A}}{x} = \sum_{x,y} \melem{y}{\op{A}}{x}\braket{x}{y} = \sum_y \melem{y}{\op{A}}{y}$; im letzten Schritt ist die Summe über $x$ die Vollständigkeitsrelation der ersten Basis. Dieselbe Rechnung zeigt die Regel, die wir ständig benutzen werden. Sind $\op{A}\colon \Hil_1 \to \Hil_2$ und $\op{B}\colon \Hil_2 \to \Hil_1$ linear, so ist
\begin{equation}
	\tr_2(\op{A}\op{B}) = \sum_{x,y} \melem{y}{\op{A}}{x}\melem{x}{\op{B}}{y} = \tr_1(\op{B}\op{A}),
\end{equation}
mit Orthonormalbasen $\{\ket{x}\}$ von $\Hil_1$ und $\{\ket{y}\}$ von $\Hil_2$; der Index an der Spur sagt, in welchem Raum sie genommen wird. Unter der Spur darf man also den letzten Faktor nach vorn holen; beliebiges Vertauschen ist nicht erlaubt. Den wichtigsten Fall liefern Ket und Bra als Abbildungen $\ket{\Psi}\colon\mathbb{C}\to\Hil$ und $\bra{\Phi}\colon\Hil\to\mathbb{C}$ aus Vorlesung~4. Mit der Basis $\{1\}$ von $\mathbb{C}$ ist die Spur der Multiplikation mit einer Zahl diese Zahl, also
\begin{equation}
	\tr\bigl(\ketbra{\Psi}{\Phi}\bigr) = \braket{\Phi}{\Psi}.
\end{equation}
Die Spur eines Ket-Bras ist ein Bra-Ket. Zum Beispiel ist der Projektor $\op{B}$ aus Vorlesung~4 gleich $\ketbra{u}{u}$ mit $\ket{u} = (1,1,1)/\sqrt{3}$, also $\tr(\op{B}) = \braket{u}{u} = 1$. In unendlich vielen Dimensionen ist die Summe eine Reihe, die nicht konvergieren muss; dort hat nicht jeder Operator eine Spur.

\subsection{Eigenwerte und Diagonalisierung}
Allgemeine Matrizen sind mühsam und vertauschen nicht miteinander; diagonale Operatoren dagegen sind so gut wie Zahlen. Deshalb muss man Operatoren oft \emph{diagonalisieren}. Alle Operatoren auf einmal diagonal zu machen gelingt in der Regel nicht; man diagonalisiert immer einen bestimmten. Den Begriff aus Vorlesung~4 machen wir jetzt genau. Er braucht kein Skalarprodukt.
\begin{definition}
	Ein Vektor $\ket{v}\neq 0$ heißt \emph{Eigenvektor} von $\op{A}$ zum \emph{Eigenwert} $\lambda\in\mathbb{C}$, wenn $\op{A}\ket{v} = \lambda\ket{v}$. Die Eigenvektoren zu $\lambda$ bilden mit dem Nullvektor den \emph{Eigenraum} zu $\lambda$.

	$\op{A}$ heißt \emph{diagonal} in der Orthonormalbasis $\{\ket{x}\}$, wenn $\melem{x}{\op{A}}{y} = \delta_{xy}\,a_x$, also $\op{A} = \sum_x a_x\ketbra{x}{x}$, und \emph{unitär diagonalisierbar}, wenn es eine Orthonormalbasis gibt, in der $\op{A}$ diagonal ist. Wir schreiben $\diag(a_1,\dots,a_n)$ für die Diagonalmatrix.
\end{definition}
Ein Eigenvektor, den man aus einem Gleichungssystem gewinnt, hat fast nie die Norm $1$; man teilt durch seine Norm und hat wieder einen Eigenvektor. Ein diagonaler Operator multipliziert jede Komponente mit ihrer eigenen Zahl, $\op{A}\ket{\Psi} = \sum_x a_x\braket{x}{\Psi}\ket{x}$; jeder Basisvektor ist also ein Eigenvektor. Das gilt auch umgekehrt, und so sind Diagonalisieren und Eigenvektoren-Suchen dieselbe Aufgabe. Sei $\{\ket{x}\}$ eine feste Orthonormalbasis und $\{\ket{\phi_x}\}$ eine Orthonormalbasis aus Eigenvektoren mit denselben Indizes, $\op{A}\ket{\phi_x} = a_x\ket{\phi_x}$. Dann ist
\begin{equation}
	\op{U} = \sum_x \ketbra{x}{\phi_x}
\end{equation}
unitär, denn mit der Orthonormalität und der Vollständigkeitsrelation ist
\begin{equation}
	\op{U}\op{U}\adj = \sum_{x,y} \ket{x}\braket{\phi_x}{\phi_y}\bra{y} = \sum_x \ketbra{x}{x} = \mathbb{I},
\end{equation}
ebenso $\op{U}\adj\op{U} = \mathbb{I}$ mit $\op{U}\adj = \sum_x\ketbra{\phi_x}{x}$, und mit $\melem{\phi_x}{\op{A}}{\phi_y} = a_y\,\delta_{xy}$ in der Mitte wird
\begin{equation}
	\op{U}\op{A}\op{U}\adj = \sum_x a_x \ketbra{x}{x}, \qquad
	\op{A} = \sum_x a_x \ketbra{\phi_x}{\phi_x}, \qquad
	\tr(\op{A}) = \sum_x a_x.
\end{equation}
Die zweite Gleichung heißt \emph{Spektralzerlegung}. Ist umgekehrt $\op{U}$ unitär und $\op{U}\op{A}\op{U}\adj$ diagonal, so sind die Vektoren $\op{U}\adj\ket{x}$ eine Orthonormalbasis aus Eigenvektoren.

In der linearen Algebra heißt ein Operator schon diagonalisierbar, wenn er irgendeine Basis aus Eigenvektoren hat; wir meinen hier stets die orthonormale, also unitäre Form. Nicht jeder Operator lässt sich diagonalisieren: Die Matrix mit den Zeilen $(0,1)$ und $(0,0)$ hat nur die Vielfachen von $(1,0)$ als Eigenvektoren. In der Quantenmechanik haben wir Glück: Die hermiteschen und die unitären Operatoren lassen sich in endlicher Dimension immer diagonalisieren. Der erste Schritt ist ein Trick, den man immer wieder braucht: denselben Ausdruck auf zwei Arten auswerten. Sei $\op{A}$ hermitesch mit $\op{A}\ket{\phi} = a\ket{\phi}$ und $\op{A}\ket{\psi} = b\ket{\psi}$. Lässt man $\op{A}$ nach rechts wirken, ist $\melem{\phi}{\op{A}}{\psi} = b\braket{\phi}{\psi}$; lässt man es nach links wirken, ist es $\braket{\op{A}\phi}{\psi} = a^*\braket{\phi}{\psi} = a\braket{\phi}{\psi}$, denn $a\braket{\phi}{\phi} = \melem{\phi}{\op{A}}{\phi}$ ist reell (Vorlesung~4). Also ist
\begin{equation}
	(a-b)\braket{\phi}{\psi} = 0:
\end{equation}
Eigenvektoren zu verschiedenen Eigenwerten sind orthogonal. Dass es immer genug davon gibt, auch wenn Eigenwerte zusammenfallen, sagt der folgende Satz; sein Beweis steht im Webkurs.
\begin{merksatz}[Spektralsatz]
	Auf einem endlichdimensionalen Hilbertraum hat jeder hermitesche und jeder unitäre Operator eine Orthonormalbasis aus Eigenvektoren, ist also unitär diagonalisierbar.
\end{merksatz}
Die endliche Dimension ist wesentlich: In unendlich vielen Dimensionen gibt es hermitesche Operatoren ganz ohne Eigenvektoren; der Webkurs zeigt ein Beispiel.

\begin{beispiel}
	Sei $\Hil = \mathbb{C}^2$, ein \emph{Qubit}, mit der Basis $\ket{0},\ket{1}$, und $\pauli{x} = \bigl(\begin{smallmatrix} 0 & 1 \\ 1 & 0\end{smallmatrix}\bigr)$. Aus $\det(\pauli{x} - \lambda\mathbb{I}) = \lambda^2 - 1$ folgen die Eigenwerte $\pm 1$ mit den normierten Eigenvektoren
	\begin{equation}
		\ket{\phi_0} = \frac{1}{\sqrt{2}}\begin{pmatrix} 1 \\ 1 \end{pmatrix}, \qquad \ket{\phi_1} = \frac{1}{\sqrt{2}}\begin{pmatrix} 1 \\ -1 \end{pmatrix}.
	\end{equation}
	Ein Ket ist eine Spalte, ein Bra eine Zeile, ein Ket-Bra also eine Matrix:
	\begin{equation}
		\op{U} = \ketbra{0}{\phi_0} + \ketbra{1}{\phi_1} = \frac{1}{\sqrt{2}}\begin{pmatrix} 1 & 1 \\ 0 & 0 \end{pmatrix} + \frac{1}{\sqrt{2}}\begin{pmatrix} 0 & 0 \\ 1 & -1 \end{pmatrix} = \frac{1}{\sqrt{2}}\begin{pmatrix} 1 & 1 \\ 1 & -1 \end{pmatrix},
	\end{equation}
	die \emph{Hadamard-Matrix}; man rechnet $\op{U}\pauli{x}\op{U}\adj = \diag(1,-1)$ nach. Vertauscht man die beiden Eigenvektoren, so liefert dasselbe Rezept eine andere unitäre Matrix, und die Diagonalform wird $\diag(-1,1)$: Die Diagonalisierung ist nicht eindeutig.
\end{beispiel}

% fig_vl05_diagonal.tex — Diagonalisierung als Basiswechsel: derselbe
% Operator, einmal in der Standardbasis, einmal in der Eigenbasis.
%
% Follows the conventions of fig_vl11_kasten.tex (the reference plate).
%
%   * geometry comes from the lecture's own worked example and nothing
%     else.  Vorlesung 5 diagonalises A = sigma_x = [[0,1],[1,0]] with the
%     Hadamard matrix U = (1/sqrt2)[[1,1],[1,-1]].  Drawn as it stands,
%     sigma_x is a reflection: it maps the unit circle onto itself and
%     there is no shape to see.  So the plate draws the shifted operator
%         Atilde = A + 2*Id = [[2,1],[1,2]],
%     which has exactly the same eigenvectors and is diagonalised by
%     exactly the same U from the lecture's Beispiel, its eigenvalues
%     merely raised by 2:
%         |phi_1> = (1,1)/sqrt2,  a_1 =  1 + 2 = 3
%         |phi_2> = (1,-1)/sqrt2, a_2 = -1 + 2 = 1
%         U Atilde U^dag = diag(3,1).
%     Script v0.3 (2026-11-04): the lecture calls sigma_x by its name and the
%     drawn operator A = sigma_x + 2 I (hat, CONV-HAT); bases |0>, |1> and
%     |phi_0>, |phi_1> as in the course figure c03.fig.diagonal.
%     Every curve below is that matrix applied to the unit circle,
%     plotted parametrically, not sketched:
%         panel (a)  t |-> (2cos t + sin t,  cos t + 2 sin t)
%         panel (b)  t |-> (3cos t,          sin t)
%     and the two are the same ellipse, seen from the two bases: U carries
%     (a) into (b).  (U is a reflection, det U = -1, not a rotation.  On
%     the centrally symmetric ellipse a -45 degree turn would give the
%     same curve, but not the same arrows: it sends |phi_2> to -e_2, where
%     U sends it to +e_2, and +e_2 is what is drawn.)  Semi-axes 3 and 1
%     are the eigenvalues, principal axes are the eigenvectors; both
%     follow from Atilde = Atilde^dag.
%     The arrow tips are the images of the basis vectors and therefore the
%     columns of the matrix: Atilde|1> = (2,1), Atilde|2> = (1,2) in (a),
%     Atilde|phi_1> = 3|phi_1>, Atilde|phi_2> = |phi_2> in (b).
%
%   * both panels are drawn at the same units-per-cm and on the same
%     ranges, because the whole claim of the figure is that the object is
%     the same and only the frame moved.  A different scale in the two
%     panels would destroy exactly that reading.
%
%   * colour encodes *which basis vector*, and nothing else: cgblue for
%     the first, birrenorange for the second.  Input and image are told
%     apart by tint and weight rather than by hue — the incoming unit
%     vector is a wide pale arrow, its image a narrower saturated one, so
%     the pair survives greyscale printing and, where the two coincide
%     (a_2 = 1 in panel (b)), the pale halo shows the coincidence instead
%     of hiding it.  The unit circle, the coordinate rays and the
%     principal axes are neutral grey: they are reference lines, not
%     data.  The image ellipse is neutral navy — it belongs to no basis,
%     which is the point.
%
%   * the dotted tip-to-tip segments in (a) say which point of the circle
%     goes to which point of the ellipse.  In (b) they are omitted: there
%     the displacement runs along the image arrow itself and a dotted
%     line on top of it would only be noise.
%
% Panel (a) standard basis.  The matrix has off-diagonal entries, the
%           ellipse lies askew, and neither basis vector comes back
%           parallel to itself.
% Panel (b) eigenbasis.  Same ellipse, axes now on the coordinate axes,
%           matrix diagonal, and the arrows only stretch: |phi_1> by the
%           factor 3, |phi_2> not at all.  Circle and ellipse touch at
%           +/- |phi_2> precisely because a_2 = 1.

\begin{center}
\begin{tikzpicture}[
    refray/.style   ={-{Stealth[length=4pt,width=3pt]}, line width=0.5pt,
                      birrengray!60},
    refcurve/.style ={line width=0.7pt, birrengray!70, densely dashed},
    princ/.style    ={line width=0.6pt, birrengray!70,
                      dash pattern=on 3.5pt off 1.5pt on 0.8pt off 1.5pt},
    imgcurve/.style ={line width=1.2pt, spacecadet},
    invec/.style    ={-{Stealth[length=6pt,width=5.5pt]}, line width=2.4pt},
    outvec/.style   ={-{Stealth[length=5.5pt,width=4pt]}, line width=1.3pt},
    movedot/.style  ={line width=0.7pt, densely dotted},
    % The four ket labels sit inside the unit circle, where the circle and
    % the principal axes leave no gap wide enough for a two-glyph label.
    % They therefore carry the same faint white backing the house legend
    % style uses, so a reference line never runs through a bracket.
    ketlbl/.style   ={lbl, font=\scriptsize, inner sep=0.9pt,
                      fill=white, fill opacity=0.85, text opacity=1},
]

% ─────────────────────────────────────────────────────────────
%  (a)  Standard basis: Atilde = [[2,1],[1,2]] is not diagonal
% ─────────────────────────────────────────────────────────────
\begin{axis}[
    name=std,
    tpcplot,
    scale only axis,
    width=5.7cm, height=5.7cm,     % equal spans below: 45 degrees looks like 45
    axis lines=none, grid=none,
    xmin=-3.35, xmax=3.35,
    ymin=-3.35, ymax=3.35,
    clip=false,
    title={(a) \tpcifenglish{Standard basis}{Standardbasis}},
]

  % ---- coordinate frame (reference) ----
  \draw[refray] (axis cs:-2.60,0) -- (axis cs:3.15,0);
  \draw[refray] (axis cs:0,-2.60) -- (axis cs:0,3.15);

  % ---- unit circle (reference): the input ----
  \addplot[refcurve, variable=t, domain=0:360, samples=181]
    ({cos(t)},{sin(t)});

  % ---- principal axes of the image ellipse = the eigen-directions,
  %      drawn only over the ellipse: semi-axes a_1 = 3, a_2 = 1 ----
  \draw[princ] (axis cs:-2.1213,-2.1213) -- (axis cs:2.1213,2.1213);
  \draw[princ] (axis cs:-0.7071,0.7071)  -- (axis cs:0.7071,-0.7071);
  \node[lbl, text=birrengray, font=\scriptsize, anchor=north east]
    at (axis cs:-2.16,-2.16) {$|\phi_0\rangle$};
  \node[lbl, text=birrengray, font=\scriptsize, anchor=north west]
    at (axis cs:0.74,-0.78) {$|\phi_1\rangle$};

  % ---- image of the unit circle under A ----
  \addplot[imgcurve, variable=t, domain=0:360, samples=181]
    ({2*cos(t)+sin(t)},{cos(t)+2*sin(t)});

  % ---- where the tips of the basis vectors go ----
  \draw[movedot, cgblue]       (axis cs:1,0) -- (axis cs:2,1);
  \draw[movedot, birrenorange] (axis cs:0,1) -- (axis cs:1,2);

  % ---- the basis vectors and their images ----
  \draw[invec,  cgblue!30]       (axis cs:0,0) -- (axis cs:1,0);
  \draw[invec,  birrenorange!30] (axis cs:0,0) -- (axis cs:0,1);
  \draw[outvec, cgblue]          (axis cs:0,0) -- (axis cs:2,1);
  \draw[outvec, birrenorange]    (axis cs:0,0) -- (axis cs:1,2);

  % |1> sits outside the ellipse to the lower right of its own arrowhead:
  % every spot nearer the origin is crossed by a ray or by the ellipse.
  \node[ketlbl, text=cgblue, anchor=north west]
    at (axis cs:1.32,-0.08) {$|0\rangle$};
  \node[ketlbl, text=birrenorange, anchor=east]
    at (axis cs:-0.12,0.70) {$|1\rangle$};
  \node[lbl, text=cgblue, font=\scriptsize, anchor=west]
    at (axis cs:2.10,0.92) {$\op{A}|0\rangle$};
  \node[lbl, text=birrenorange, font=\scriptsize, anchor=south]
    at (axis cs:1.00,2.14) {$\op{A}|1\rangle$};

  % ---- the matrix this frame produces ----
  % Two lines: on one line the definition and the matrix together reach
  % as far as the note in the opposite corner.
  \node[lbl, anchor=north west, align=left]
    at (axis cs:-3.30,3.22)
    {$\op{A} = \pauli{x} + 2\mathbb{I}$\\[3pt]
     $\phantom{\op{A}} = \begin{pmatrix} 2 & 1 \\ 1 & 2\end{pmatrix}$};

  \node[lbl, text=birrengray, font=\scriptsize, anchor=north east]
    at (axis cs:3.30,3.18)
    {\tpcifenglish{circle $\mapsto$ ellipse}{Kreis $\mapsto$ Ellipse}};

  \node[lbl, text=birrengray, font=\scriptsize, anchor=north]
    at (axis cs:0,-2.74)
    {\tpcifenglish{no basis vector is an eigenvector}%
                  {kein Basisvektor ist Eigenvektor}};
\end{axis}

% ─────────────────────────────────────────────────────────────
%  (b)  Eigenbasis: same ellipse, U Atilde U^dag = diag(3,1)
% ─────────────────────────────────────────────────────────────
\begin{axis}[
    tpcplot,
    scale only axis,
    at={(std.south east)}, anchor=south west, xshift=20mm,
    width=5.7cm, height=5.7cm,
    axis lines=none, grid=none,
    xmin=-3.35, xmax=3.35,
    ymin=-3.35, ymax=3.35,
    clip=false,
    title={(b) \tpcifenglish{Eigenbasis}{Eigenbasis}},
]

  % ---- coordinate frame: now the eigen-directions ----
  \draw[refray] (axis cs:-2.60,0) -- (axis cs:3.25,0);
  \draw[refray] (axis cs:0,-2.60) -- (axis cs:0,3.15);

  % ---- unit circle (reference) ----
  \addplot[refcurve, variable=t, domain=0:360, samples=181]
    ({cos(t)},{sin(t)});

  % ---- the same ellipse, now axis-aligned: semi-axes a_1 = 3, a_2 = 1 ----
  \addplot[imgcurve, variable=t, domain=0:360, samples=181]
    ({3*cos(t)},{sin(t)});

  % ---- the basis vectors and their images: pure stretching ----
  \draw[invec,  cgblue!30]       (axis cs:0,0) -- (axis cs:1,0);
  \draw[invec,  birrenorange!30] (axis cs:0,0) -- (axis cs:0,1);
  \draw[outvec, cgblue]          (axis cs:0,0) -- (axis cs:3,0);
  \draw[outvec, birrenorange]    (axis cs:0,0) -- (axis cs:0,1);

  % Both labels are kept inside the unit circle: at their old places the
  % dashed circle ran straight through the ket brackets.
  \node[ketlbl, text=cgblue, anchor=north]
    at (axis cs:0.45,-0.14) {$|\phi_0\rangle$};
  \node[ketlbl, text=birrenorange, anchor=east]
    at (axis cs:-0.10,0.58) {$|\phi_1\rangle$};

  \node[lbl, text=cgblue, font=\scriptsize, anchor=north]
    at (axis cs:1.75,-1.10) {$\op{A}|\phi_0\rangle = 3\,|\phi_0\rangle$};
  \node[lbl, text=birrenorange, font=\scriptsize, anchor=south west]
    at (axis cs:0.14,1.10) {$\op{A}|\phi_1\rangle = |\phi_1\rangle$};

  % ---- the matrix this frame produces ----
  \node[lbl, anchor=north west] at (axis cs:-3.30,3.22)
    {$\op{U}\op{A}\op{U}^\dagger
       = \begin{pmatrix} 3 & 0 \\ 0 & 1\end{pmatrix}$};

  \node[lbl, text=birrengray, font=\scriptsize, anchor=north]
    at (axis cs:0,-2.74)
    {\tpcifenglish{every basis vector is an eigenvector}%
                  {jeder Basisvektor ist Eigenvektor}};
\end{axis}

% ─────────────────────────────────────────────────────────────
%  The basis change between the two frames.  U is the Hadamard matrix
%  the lecture's Beispiel already computed.
% ─────────────────────────────────────────────────────────────
\draw[-{Stealth[length=5pt,width=3.5pt]}, spacecadet, line width=0.8pt]
  ([xshift=6mm,yshift=8mm]std.east) -- ([xshift=14mm,yshift=8mm]std.east);
\node[lbl, font=\scriptsize, anchor=south]
  at ([xshift=10mm,yshift=8.6mm]std.east) {$\op{U}$};

\end{tikzpicture}
\end{center}


Die Abbildung zeigt dieses Beispiel, aber nicht $\pauli{x}$ selbst: Als Spiegelung bildet $\pauli{x}$ den Einheitskreis auf sich ab. Gezeichnet ist $\op{A} = \pauli{x} + 2\mathbb{I}$, mit denselben Eigenvektoren und derselben Matrix $\op{U}$ und den Eigenwerten $3$ und $1$. Gezeichnet ist das Bild des Einheitskreises von $\mathbb{R}^2$. Links liegt die Bildellipse schief, rechts, in der Eigenbasis, liegt sie auf den Achsen. Ihre Hauptachsen liegen auf den Eigenvektoren, ihre Halbachsen sind die hier positiven Eigenwerte.

\subsection{Funktionalkalkül}
Weil diagonale Operatoren so gut wie Zahlen sind, kann man Funktionen auf sie anwenden, indem man die Funktion auf die Zahlen anwendet.
\begin{definition}
	Sei $\op{A} = \sum_x a_x\ketbra{\phi_x}{\phi_x}$ eine Spektralzerlegung und $f$ eine komplexe Funktion auf den Eigenwerten. Dann ist $f(\op{A}) = \sum_x f(a_x)\ketbra{\phi_x}{\phi_x}$.
\end{definition}
Für ein Polynom stimmt das mit dem Einsetzen überein, denn $\op{A}^k = \sum_x a_x^k\ketbra{\phi_x}{\phi_x}$; und weil jede Funktion auf den endlich vielen Eigenwerten mit einem Polynom übereinstimmt (Lagrange-Interpolation, im Webkurs), hängt $f(\op{A})$ nicht von der gewählten Eigenbasis ab. Der wichtigste Fall: Sei $\op{\Ham} = \sum_x E_x\ketbra{\phi_x}{\phi_x}$ hermitesch, in Vorlesung~6 der Hamiltonoperator mit den Energien $E_x$, und $t$ reell. Dann ist
\begin{equation}
	\op{U}_t = e^{-\frac{it}{\hbar}\op{\Ham}} = \sum_x e^{-\frac{it}{\hbar}E_x}\ketbra{\phi_x}{\phi_x}.
\end{equation}
Auf jedem Eigenvektor wirkt $\op{U}_t$ als Multiplikation mit einer \emph{Phase}. Die Kette ist: Die $E_x$ sind reell, weil $\op{\Ham}$ hermitesch ist; deshalb haben die Phasen den Betrag $1$; deshalb ist $\op{U}_t\adj\op{U}_t = \sum_x \abs{e^{-itE_x/\hbar}}^2\ketbra{\phi_x}{\phi_x} = \mathbb{I}$, und $\op{U}_t$ ändert die Norm keines Vektors. In endlicher Dimension ist $\op{U}_t$ auch die Summe der Exponentialreihe $\sum_n \frac{1}{n!}\bigl(\frac{-it}{\hbar}\bigr)^n\op{\Ham}^n$, denn jede Teilsumme ist ein Polynom in $\op{\Ham}$, wirkt also auf $\ket{\phi_x}$ als die Teilsumme der Zahlenreihe für $e^{-itE_x/\hbar}$, und diese endlich vielen Reihen konvergieren.

\begin{beispiel}
	Sei $\op{\Ham} = \varepsilon\pauli{x}$ mit einer Energie $\varepsilon>0$. Mit den Eigenwerten $\pm\varepsilon$ und den Eigenvektoren von oben ist
	\begin{equation}
		\eqbox{\op{U}_t = e^{-i\varepsilon t/\hbar}\ketbra{\phi_0}{\phi_0} + e^{i\varepsilon t/\hbar}\ketbra{\phi_1}{\phi_1} = \cos\Bigl(\frac{\varepsilon t}{\hbar}\Bigr)\mathbb{I} - i\sin\Bigl(\frac{\varepsilon t}{\hbar}\Bigr)\pauli{x}}
	\end{equation}
	wegen $\mathbb{I} = \ketbra{\phi_0}{\phi_0} + \ketbra{\phi_1}{\phi_1}$ und $\pauli{x} = \ketbra{\phi_0}{\phi_0} - \ketbra{\phi_1}{\phi_1}$. Dasselbe liefert die Reihe: Wegen $\pauli{x}^2 = \mathbb{I}$ zerfällt sie in die Reihen für Kosinus und Sinus. Dieser Trick nutzt $\pauli{x}^2 = \mathbb{I}$; im Allgemeinen rechnet man mit der Spektralzerlegung.
\end{beispiel}

\begin{merksatz}
	Der Operator $\op{U}_t$ beschreibt die Zeitentwicklung eines Quantensystems. Damit steht die Lösung der Schrödingergleichung schon hier, bevor die Gleichung selbst aufgetaucht ist. „Die Dynamik lösen“ heißt deshalb: den Hamiltonoperator diagonalisieren.
\end{merksatz}
